% !TeX program = lualatex
% !TeX encoding = utf8
% !TeX spellcheck = uk_UA

\documentclass{article}
\usepackage[fontsize=12pt]{fontsize}

\usepackage{fontsetup}
\usepackage[english, russian, ukrainian]{babel}
\usepackage{microtype}
\usepackage{tabularray}
\usepackage{xcolor}


\usepackage[%
	a4paper,%
	footskip=1cm,%
	headsep=0.3cm,%
	top=2cm, %поле сверху
	bottom=2cm, %поле снизу
	left=2cm, %поле ліворуч
	right=2cm, %поле праворуч
    ]{geometry}

\usepackage{amsmath}
\usepackage{graphicx}
\usepackage{floatflt}

\usepackage[%colorlinks=true,
	%urlcolor = blue, %Colour for external hyperlinks
	%linkcolor  = malina, %Colour of internal links
	%citecolor  = green, %Colour of citations
	bookmarks = true,
	bookmarksnumbered=true,
	unicode,
	linktoc = all,
	hypertexnames=false,
	pdftoolbar=false,
	pdfpagelayout=TwoPageRight,
	pdfauthor={Ponomarenko S.M. aka sergiokapone},
	pdfdisplaydoctitle=true,
	pdfencoding=auto
	]%
	{hyperref}



\title{Рейтингова система оцінювання (РСО) \\(витримки з \href{https://osvita.kpi.ua/sites/default/files/downloads/Pologennia_RSO_2025.pdf}{документу})}
\date{}

\begin{document}
\maketitle

\section{Типи рейтингових систем}

Згідно з «Положенням про систему оцінювання результатів навчання в КПІ ім. Ігоря Сікорського» (ред. 2025):

\begin{itemize}
    \item \textbf{РСО-1} --- застосовується для дисциплін, семестровий контроль з яких проводиться у формі \textbf{заліку}.
    Вона містить лише поточний контроль і не має окремої екзаменаційної складової.
    \item \textbf{РСО-2} --- застосовується для дисциплін, семестровий контроль з яких проводиться у формі \textbf{екзамену},
    а також для курсових проєктів, практик тощо. Складається з \textbf{двох складових}:
    стартової (поточний контроль) і екзаменаційної.
    \item Стартова складова --- це «накопичувальна частина»: наскільки студент працював протягом семестру.
\end{itemize}



\subsection*{РСО-1 (залік)}

\noindent\textbf{Тип системи:} РСО-1 \\
\textbf{Форма семестрового контролю:} залік

\begin{table}[h!]
\centering
\caption{Розподіл балів за видами контролю (РСО-1, залік)}
\begin{tblr}{
colspec = {Q[l,m, 7cm]|X|Q[c]|},
row{1}={c,m, font=\bfseries, bg=cyan!20},
hlines, vlines
}
{Вид контролю / діяльності} & {Форма проведення} & {Макс. бали} \\
Модульна контрольна робота & Синхронна & 40 \\
Домашня контрольна робота / індивідуальне семестрове завдання & Асинхронна & 40 \\
Поточна активність, короткі тести, відвідування, участь у дискусіях & Синхронна / асинхронна & 10 \\
Залік (для підвищення рейтингу або у разі $<60$ балів) & Семестровий контроль & 10 \\
\SetCell[r=1,c=2]{r}\textbf{Загальна кількість балів} & & \textbf{100} \\ \end{tblr}
\end{table}

\noindent\textit{Пояснення:} студент, який виконав усі умови поточного контролю та набрав $\geq 60$ балів, отримує залік автоматично.
При бажанні може пройти залік для підвищення рейтингу.
Тип оцінювання --- \textbf{РСО-1 (лише поточна складова)}.


\section{РСО-2 (екзамен)}

\noindent\textbf{Тип системи:} РСО-2 \\
\textbf{Форма семестрового контролю:} екзамен

\begin{table}[h!]
\centering
\caption{Розподіл балів за видами контролю (РСО-2, екзамен)}
\begin{tblr}{
colspec = {Q[l,m, 7cm]|X|Q[c]|},
row{1}={c,m, font=\bfseries, bg=cyan!20},
hlines, vlines
}
{Вид контролю / діяльності} & {Форма проведення} & {Макс. бали} \\
Модульна контрольна робота & Синхронна & 20 \\
Розрахункова робота / лабораторні / практичні & Асинхронна / синхронна & 20 \\
Поточна активність, короткі тести & Синхронна / асинхронна & 10 \\
\SetCell[r=1,c=2]{r}\textbf{Разом за поточний контроль (стартова складова)} &  & \textbf{50} \\
Екзамен (письмова або усна форма згідно з силабусом) & Семестровий контроль & 50 \\
\SetCell[r=1,c=2]{r}\textbf{Загальна кількість балів} &  & \textbf{100} \\
\end{tblr}
\end{table}

\noindent\textit{Мінімальна кількість балів для допуску до екзамену:} 25–30 (50–60\% від стартової складової).
\textit{Фінальна рейтингова оцінка:}
\[
R_{\text{підсумк.}} = R_{\text{старт.}} + R_{\text{екз.}}, \quad 0 \le R_{\text{підсумк.}} \le 100.
\]
Тип оцінювання — \textbf{РСО-2 (стартова + екзаменаційна складові)}.


\subsection*{Шкала оцінювання КПІ (Додаток А)}

\begin{table}[h!]
\centering
\renewcommand{\arraystretch}{1.2}
\begin{tabular}{|c|c|c|}
\hline
\textbf{Бали} & \textbf{Оцінка} & \textbf{Рівень} \\ \hline
95--100 & Відмінно & A \\ \hline
85--94 & Дуже добре & B \\ \hline
75--84 & Добре & C \\ \hline
65--74 & Задовільно & D \\ \hline
60--64 & Достатньо & E \\ \hline
<60 & Незадовільно & F \\ \hline
\end{tabular}
\end{table}




\end{document}


